\documentclass[11pt]{article}
\usepackage[margin=1in]{geometry}
\usepackage{booktabs}
\usepackage{amsmath,amssymb}
\usepackage{threeparttable}
\usepackage[hidelinks]{hyperref}
\usepackage{ctex}

\title{连续行业专长的 GPS-IPW 加权回归报告}
\author{}
\date{}

\begin{document}
\maketitle

\section{研究目的与方法定位}

本文将事务所行业专长市场份额 $IMS$ 作为连续处理变量，使用连续处理的广义倾向得分（generalized propensity score, GPS）构造稳定化逆条件密度权重。研究目的不是将连续 $IMS$ 任意二元化，而是在处理前企业特征更具可比性的伪总体中，考察不同程度的行业专长与企业漂绿之间的条件关联。

连续处理 GPS 的基本思想是对每一个处理水平 $t$ 估计条件处理密度
\[
r(t,x)=f_{T\mid X}(t\mid x),
\]
并利用该密度建立处理前协变量的平衡。Imbens (2000) 证明了连续处理情形下 GPS 的平衡性质；Hirano and Imbens (2004) 给出了基于条件处理分布和剂量--反应函数的实施框架。经济管理领域的应用也遵循相同原则：Bia and Mattei (2012) 根据企业获得不同金额财政援助的制度分配机制选择处理前企业特征；Kluve et al. (2012) 纳入培训分配者实际使用的教育、就业历史和地区劳动力市场信息；Ammar and Boughrara (2023) 则根据银行多元化决策机制纳入银行特征以及宏观、行业环境变量。

本报告中的 GPS-IPW 是一种基于处理前可观测变量的选择调整方法。它不能自动解决未观测的企业--审计师匹配因素、反向因果或行业内溢出，因此结果主要解释为条件关联和稳健性证据，而不是无条件的因果证明。

\section{数据、缩尾口径与分析样本}

主回归在连续控制变量上采用 1\% 和 99\% 缩尾；$IMS$、Greenwashing 以及其他行业专长指标不缩尾。GPS 分析直接读取与主回归变量定义一致的共同样本文件 \texttt{Regression/outputs/paper\_tables/common\_sample\_14778.csv}，不在 GPS 脚本中再次对原始变量进行第二次缩尾，以避免 GPS 和主回归使用不同的变量构造或缩尾边界。

数据按企业代码和年份排序。只有企业内前一观测年份与当期严格相差 1 年时，才构造滞后变量。因此，GPS 第一阶段使用真正的 $t-1$ 信息，而不会将非连续年份错误地当作相邻年度。由于第一阶段需要当期 $IMS$、当期 Greenwashing 以及完整的滞后共同原因变量，最终 GPS 共同分析样本为 10,001 个公司年度观测，涉及 4,007 家企业。该样本量低于主回归的 14,778 个公司年度观测，主要源于滞后变量和连续年度要求，不应将其解释为 GPS 权重造成的样本损失。

结果变量 $GW_{i,t}$ 为 \texttt{Greenwashing连续\_基准}。连续处理变量 $IMS_{i,t}$ 为事务所 $a$ 在企业所属行业--年度市场中的客户营业收入份额。$IMS$ 保留其原始经济含义，不进行 1\%--99\% 变量缩尾；仅在后续权重处理中考察极端权重的 P0.5--P99.5 和 P1--P99 截尾稳健性。

\section{第一阶段协变量的选择逻辑}

\subsection{协变量选择的因果原则}

第一阶段的任务是估计企业在处理发生前的可观测信息条件下进入不同 $IMS$ 水平的条件密度，而不是建立一个以拟合优度为目标的预测模型。本文遵循弱无混杂假设：对于任一可行处理水平 $t$，在给定处理前协变量 $X_{i,t-1}$ 后，企业接受的 $IMS_{i,t}$ 与其潜在 Greenwashing 结果独立。因而，协变量是否进入第一阶段，取决于其是否在 $IMS$ 形成之前已经确定，且可能同时影响企业匹配到不同程度行业专长审计师的概率与其 Greenwashing；或者是否是企业和审计师在匹配决策中能够观察、使用的重要信息。

这一做法与 Bia and Mattei (2012) 的思路一致。该文并未按照处理模型中单个变量的显著性、$p$ 值、调整后 $R^2$ 或逐步回归结果自动筛选协变量，而是结合财政援助的制度分配背景，将其表 2 和表 3 所列的处理前企业特征整体纳入 GPS 估计，并在估计后通过均值差异和 Bayes Factor 评价协变量平衡。本文同样以企业--审计师匹配和行业专长形成机制为基础预先设定协变量集合；SMD、条件密度、ESS 和权重分布仅用于评价所设定 GPS 的可比性和稳定性，不能反过来作为某一个协变量进入模型的唯一理论依据。

因此，不应因某个第一阶段系数不显著而删除具有明确混杂路径含义的变量。第一阶段变量选择服务于弱无混杂假设，而非单纯追求 $IMS$ 的预测精度。同理，$|SMD|<0.1$、Bayes Factor、$p$ 值、调整后 $R^2$ 和 ESS 都是估计后的诊断信息，不是事前变量筛选规则。

\subsection{主规格设定}

本文在正文中采用\textbf{完整共同原因加审计选择规格}作为主 GPS 第一阶段规格。该选择的首要依据是：$IMS$ 反映企业与不同程度行业专长审计师的匹配结果，处理分配过程不仅取决于企业既有的财务、治理和披露特征，也直接受到既有审计师类型、审计投入、审计关系持续时间及换所决策的影响。若忽略后一组处理前审计关系信息，条件密度模型可能未能充分刻画企业进入不同 $IMS$ 水平的可观测选择渠道。

令 $Z_{i,t}=logit(IMS_{i,t})$，主规格为：
\[
Z_{i,t}=\eta'X_{i,t-1}+\lambda_t+\delta_j+u_{i,t},
\qquad u_{i,t}\sim N(0,\sigma^2),
\]
其中，
\[
\begin{split}
X_{i,t-1}=\{&L\_Greenwashing, L\_Size, L\_Lev, L\_ROA, L\_CFO, L\_Growth, L\_FirmAge, L\_SOE,\\
&L\_Dual, L\_BoardSize, L\_IndependentDirectorRatio, L\_Big4, L\_LnAuditFee,\\
&L\_AuditorTenure, L\_IsChangeFirm\}.
\end{split}
\]
$\lambda_t$ 为年度固定效应，$\delta_j$ 为行业固定效应。所有协变量均为严格滞后值，且仅在企业相邻年份连续时构造；这一处理保证协变量在当期 $IMS$ 和当期 Greenwashing 之前已确定，从而缓解同时性和反向因果担忧，但不能消除未观测混杂。

主规格变量具有如下因果和经济含义。滞后 Greenwashing 刻画企业既有环境沟通--绩效缺口、披露倾向、声誉压力和治理基础；这些特征既可能提高企业选择行业专长审计师的需求，也可能延续至当期 Greenwashing。Size 反映业务复杂度、审计需求与支付高质量审计服务的能力，并可能影响外部关注和漂绿空间。Lev、ROA 与 CFO 分别代理融资约束和违约风险、经营表现、内部现金流与资源约束；它们会影响客户风险、审计师承接和配置决策，也可能影响管理层使用环境叙事缓冲融资或业绩压力的动机。Growth 与 FirmAge 反映企业成长阶段、生命周期和组织成熟度，因而同时关联审计复杂度、行业专长需求及信息披露策略。SOE 反映产权性质及其监管、政治联系和外部监督环境，可能影响审计师选择和 Greenwashing 的约束程度。

滞后 Dual、BoardSize 和 IndependentDirectorRatio 刻画内部治理与监督需求。两职合一可能弱化董事会监督并影响管理层对审计质量的需求；董事会规模和独立董事比例反映监督资源、专业意见和对外部保证的需求。上述治理安排既可能影响企业匹配行业专长审计师的意愿，也可能影响 Greenwashing 的内部约束。

滞后 Big4、LnAuditFee、AuditorTenure 和 IsChangeFirm 则直接代理企业--审计师匹配机制。Big4 反映既有审计师的品牌、资源和客户筛选标准；审计费用反映业务复杂度、审计投入和客户风险；审计任期反映关系专用知识与独立性之间的动态权衡；换所变量反映企业和审计师关系的重新匹配。它们均可能影响企业在当期匹配到高行业专长审计师的机会，也可能通过审计监督强度、报告质量和管理层机会主义空间影响 Greenwashing。其纳入前提是变量均使用严格滞后值，因而被视为处理发生前的匹配信息，而非当期 $IMS$ 的后果或中介变量；这一时间顺序假设应在实证解释中保持审慎。

年度固定效应吸收年度审计市场供给、监管政策、宏观冲击和总体环境披露趋势；行业固定效应吸收行业竞争、行业风险、行业审计市场结构和行业层面的专长分配规则。由于 $IMS$ 是行业--年度中的相对市场份额，这两类固定效应是处理分配环境的组成部分，而非仅为提高回归拟合度而加入。

第一阶段与 G2 使用重叠协变量不构成重复控制：前者用于构造协变量与连续处理更可比的加权伪总体，后者用于在该伪总体中进一步调整结果方程。G2 因而应表述为 GPS 加权后的增强结果方程，而不应在未建立相应理论条件和推断程序时称为严格的双重稳健估计。

\subsection{规格选择和诊断依据}

三组规格的定位明确区分。核心共同原因规格仅包含既有 Greenwashing、财务、生命周期和产权信息；完整共同原因规格在此基础上加入治理变量；完整共同原因加审计选择规格进一步加入处理前审计师关系与换所信息。正文采用第三组作为主规格，是因为其对企业进入不同 $IMS$ 水平的选择过程具有最完整的经济解释，而非因为其调整后 $R^2$ 最高或回归结果更显著。核心共同原因和完整共同原因规格作为稳健性检验报告，以检验结论是否依赖于治理和审计选择信息的纳入。

三组规格均使用 10,001 个公司年度观测和 4,007 家企业。核心、完整和主规格的处理模型调整后 $R^2$ 分别为 0.1896、0.1900 和 0.1995；相应残差标准差分别为 1.4480、1.4477 和 1.4391。主规格的拟合度提高仅说明审计选择变量解释了部分可观测的 IMS 分配差异，并不是其被选为主规格的决定性证据。

平衡诊断提供辅助评价。GPS 加权后，核心、完整和主规格中满足 $|SMD|<0.1$ 的协变量--剂量组组合比例分别为 94.55\%、92.50\% 和 93.33\%，平均绝对 SMD 分别为 0.0406、0.0441 和 0.0426。主规格的最大绝对 SMD 仍为约 0.1666，表明平衡明显改善但并非完全。权重诊断亦显示，加入审计选择变量会改变条件密度和 ESS；主规格的未截尾、P0.5--P99.5 截尾和 P1--P99 截尾 ESS 分别为 7,197.35、8,078.75 和 8,437.53。因而，主规格并不在每一项诊断上机械占优，而是在更完整的处理分配解释、可接受的加权平衡和没有样本共同支持损失之间取得审慎权衡。

三组规格在未截尾和权重截尾下的 $IMS$ 系数方向总体为负，尽管系数大小和统计精度有所变化。这一稳定性支持负向条件关联并非仅由某一狭窄第一阶段设定驱动；但它不证明不存在未观测的企业--审计师匹配因素。

\subsection{GPS 平衡性检验}

为检验 GPS 加权是否改善可比性，本文将连续 $IMS$ 按五分位划分为五个剂量组，并分别比较每个剂量组与其余样本在各滞后协变量上的未加权与 GPS 加权均值及标准化均值差异（SMD）。剂量组仅是平衡诊断工具，不改变 $IMS$ 作为连续处理变量的设定，也不用于替代连续剂量的后续估计。

对于剂量组 $g$ 和协变量 $X$，SMD 表示为：
\[
SMD_g=\frac{\bar X_g-\bar X_{-g}}{\sqrt{(s_g^2+s_{-g}^2)/2}}.
\]
加权诊断以稳定化 GPS 权重计算相应的均值和方差。$|SMD|<0.1$ 被用作经验性平衡参考阈值，用于描述调整后可比性的改善程度，而非决定某变量是否应进入第一阶段的筛选门槛。当前主规格的加权后 $|SMD|<0.1$ 比例为 93.33\%、平均绝对 SMD 为 0.0426，但最大绝对 SMD 约为 0.1666。因此，应将结果表述为 GPS 加权显著改善了处理前可观测协变量的平衡，而不能表述为所有协变量已经完全平衡。

此外，平衡改善不能检验、更不能证明弱无混杂假设成立。GPS 识别仍依赖于：在给定上述处理前协变量、年度和行业环境后，不同 $IMS$ 水平的分配与潜在 Greenwashing 结果独立。未观测的管理者偏好、审计师私有信息或行业内溢出仍可能违背该假设。因此，GPS-IPW 结果应作为对可观测选择偏差进行调整后的条件关联和稳健性证据，与其他识别设计和敏感性分析共同解读。

\section{GPS-IPW 的计算}

由于 $IMS\in(0,1)$，先进行 logit 变换：
\[
Z_{i,t}=\log\left(\frac{IMS_{i,t}}{1-IMS_{i,t}}\right).
\]
对每个第一阶段规格分别估计：
\[
Z_{i,t}=\eta'X_{i,t-1}+\lambda_t+\delta_j+u_{i,t},
\qquad u_{i,t}\sim N(0,\sigma^2),
\]
其中 $X_{i,t-1}$ 为对应规格中的全部滞后协变量，$\lambda_t$ 和 $\delta_j$ 分别为年度和行业固定效应。条件密度为：
\[
\widehat f_{Z\mid X}(Z_{i,t}\mid X_{i,t-1})
=\phi\left(Z_{i,t};\widehat\mu_{i,t},\widehat\sigma\right),
\]
其中 $\widehat\mu_{i,t}$ 为处理模型拟合值。

稳定化 GPS-IPW 定义为：
\[
w_{i,t}=\frac{\widehat f_Z(Z_{i,t})}
{\widehat f_{Z\mid X}(Z_{i,t}\mid X_{i,t-1})}.
\]
分子为 $Z=logit(IMS)$ 的边际正态密度，分母为给定处理前协变量、年份和行业后的条件密度。权重计算对每个第一阶段规格独立进行。有效样本量采用：
\[
ESS=\frac{(\sum_iw_i)^2}{\sum_iw_i^2}.
\]

脚本报告三种权重处理：未截尾、P0.5--P99.5 截尾和 P1--P99 截尾。后两种是对权重而非原始控制变量或 $IMS$ 的稳定性处理，目的是评价少数低条件密度观测对结果的影响。

\section{第一阶段拟合与权重诊断}

\begin{table}[htbp]
\centering
\begin{threeparttable}
\caption{第一阶段条件处理模型诊断}
\begin{tabular}{lccc}
\toprule
 & 核心共同原因 & 完整共同原因 & 完整共同原因加审计选择 \\
\midrule
样本数 & 10,001 & 10,001 & 10,001 \\
企业数 & 4,007 & 4,007 & 4,007 \\
处理模型残差标准差 & 1.4480 & 1.4477 & 1.4391 \\
调整后 $R^2$ & 0.1896 & 0.1900 & 0.1995 \\
$IMS$ P5 & 0.0016 & 0.0016 & 0.0016 \\
$IMS$ P95 & 0.2098 & 0.2098 & 0.2098 \\
条件密度 P01 & 0.0061 & 0.0062 & 0.0060 \\
条件密度 P50 & 0.2175 & 0.2178 & 0.2201 \\
条件密度 P99 & 0.2755 & 0.2756 & 0.2772 \\
\bottomrule
\end{tabular}
\begin{tablenotes}\footnotesize
\item 注：三个规格使用同一 GPS 共同分析样本。完整共同原因规格不包含审计选择变量；完整共同原因加审计选择规格在此基础上加入 Big4、审计收费、审计任期和换所状态。
\end{tablenotes}
\end{threeparttable}
\end{table}

加入治理变量后，处理模型调整后 $R^2$ 从 0.1896 小幅提高到 0.1900，残差标准差为 1.4477；进一步加入审计选择变量后，调整后 $R^2$ 提高到 0.1995，残差标准差降至 1.4391。这说明审计选择变量能够解释一部分企业获得不同 IMS 强度的差异，但拟合度本身不是选择第一阶段规格的唯一依据。更重要的判断是加入变量后是否改善处理前协变量平衡，以及是否造成不可接受的共同支持损失。

\section{IMS 分布与剂量分组}

为实施连续处理的平衡诊断，脚本将共同样本中的 $IMS$ 按五分位划分为五个剂量组，并输出完整的分位点和剂量组分布。五个实际剂量区间约为：
\begin{center}
\begin{tabular}{lc}
\toprule
剂量组 & $IMS$ 区间 \\
\midrule
第 1 组 & $(-0.00099, 0.01015]$ \\
第 2 组 & $(0.01015, 0.02858]$ \\
第 3 组 & $(0.02858, 0.06161]$ \\
第 4 组 & $(0.06161, 0.11851]$ \\
第 5 组 & $(0.11851, 0.73073]$ \\
\bottomrule
\end{tabular}
\end{center}

这里第一个区间的左端点是五分位分组算法产生的显示边界，并不意味着 $IMS$ 取负值；$IMS$ 的实际下界仍由原始数据决定。由于连续处理的剂量组只是平衡诊断工具，而不是将 $IMS$ 重新定义为离散处理，后续回归仍使用原始连续 $IMS$。

\section{GPS 协变量平衡诊断}

\subsection{诊断方法}

连续处理 GPS 中的“匹配”不是将每一家企业机械地一对一配对，也不是只比较二元处理组和控制组。其含义是：在给定处理水平的 GPS 后，具有相近处理分配密度的企业应当在处理前协变量上更加相似。因此，脚本在五个 $IMS$ 剂量组内，分别比较该剂量组与其余样本在每一个第一阶段滞后协变量上的均值和标准化均值差异（standardized mean difference, SMD）。

对某一个剂量组 $g$ 和协变量 $X$，未调整 SMD 计算为：
\[
SMD_g=\frac{\bar X_g-\bar X_{-g}}
{\sqrt{(s_g^2+s_{-g}^2)/2}}.
\]
GPS 调整后的版本使用原始稳定化 GPS 权重计算加权均值和加权方差。脚本将每个第一阶段规格、每个 IMS 剂量组、每个协变量以及“未调整/GPS 原始权重”分别写入 \texttt{gps\_covariate\_balance.csv}，同时记录 SMD 和 $|SMD|<0.1$ 指标。

因此，当前报告已经包含实际运行的逐剂量、逐协变量平衡输出，而不是仅提出未来应当进行平衡检验。需要注意，脚本目前将“该剂量组”与“其余样本”进行比较，并未进一步在每个剂量组内部按 GPS 中位数建立子块。因此，该诊断是基于 GPS 权重的剂量组平衡检查，不应表述为完整的 GPS 分块匹配检验。

\subsection{平衡结果}

根据 \texttt{gps\_covariate\_balance.csv}，三组第一阶段规格的总体结果如下：
\begin{table}[htbp]
\centering
\begin{tabular}{lcc}
\toprule
第一阶段规格 & 未调整时 $|SMD|<0.1$ 比例 & GPS 加权后 $|SMD|<0.1$ 比例 \\
\midrule
核心共同原因 & 81.82\% & 94.55\% \\
完整共同原因 & 75.00\% & 92.50\% \\
完整共同原因加审计选择 & 74.67\% & 93.33\% \\
\bottomrule
\end{tabular}
\end{table}

总体而言，GPS 加权后满足 $|SMD|<0.1$ 的协变量--剂量组组合比例由未调整时的 74.67--81.82\% 提高至 92.50--94.55\%。以完整共同原因加审计选择规格为例，平均绝对 SMD 从 0.0821 降至 0.0426；核心共同原因规格从 0.0715 降至 0.0406；完整共同原因规格从 0.0840 降至 0.0441。这表明 GPS 权重总体上明显降低了不同 IMS 剂量组之间的处理前协变量差异。

但平衡并非完全实现。GPS 加权后的最大绝对 SMD 约为 0.1666，仍有少数组合超过 0.10。因此，当前证据应表述为“GPS 加权显著改善了可观测协变量平衡”，而不能表述为“所有协变量已经完全平衡”。后续若将 GPS 作为正式因果识别依据，应进一步检查这些残余不平衡是否集中于特定剂量组或审计选择变量，并考虑扩大条件密度模型的非线性项或采用更灵活的 GPS 估计方法。

\subsection{共同支持与极端权重}

$IMS$ 的分组范围覆盖从接近零到高处理强度的观测；高剂量组的上界约为 0.95。处理模型条件密度的 P01 在完整共同原因加审计选择规格下为 0.0060，表明需要关注条件密度较低的观测。权重诊断显示，核心共同原因规格的未截尾、P0.5--P99.5 和 P1--P99 截尾 ESS 分别为 7,001.86、8,390.22 和 8,611.55；完整共同原因规格分别为 6,856.92、8,345.77 和 8,596.82；完整共同原因加审计选择规格分别为 7,197.35、8,078.75 和 8,437.53。

权重截尾提高 ESS，说明少数极端权重对估计精度有影响；但截尾也改变了目标伪总体，因此未截尾和截尾结果应同时报告。不能因为某一种截尾规则得到更显著的 G2 就选择性报告该规则。

\section{加权结果方程}

\subsection{G1：纯 GPS-IPW 边际模型}

G1 不在结果方程中重复加入第一阶段协变量：
\[
GW_{i,t}=\alpha+\beta IMS_{i,t}+\lambda_t+\delta_j+\varepsilon_{i,t}.
\]

\subsection{G2：增强 GPS-IPW 模型}

G2 在相同第一阶段权重下加入全部滞后共同原因变量：
\[
GW_{i,t}=\alpha+\beta IMS_{i,t}+\gamma'X_{i,t-1}+\lambda_t+\delta_j+\varepsilon_{i,t}.
\]

G2 是 GPS 加权后的增强结果方程。它不应被无条件称为严格意义上的双重稳健估计，因为当前实现没有建立相应的双重稳健估计量和推断程序。G2 的系数变化表示，在调整处理分配和结果方程之后，$IMS$ 与漂绿之间剩余的条件关联发生了变化。

\section{主要回归结果}

下表报告三组第一阶段协变量规格在三种权重处理下的 G1/G2 结果。三组规格的处理模型和权重均独立估计，因此其差异反映治理变量与审计选择变量进入处理分配模型后对目标伪总体和有效信息量的影响。

\begin{table}[htbp]
\centering
\begin{threeparttable}
\caption{不同第一阶段协变量和权重处理下的 GPS-IPW 结果}
\scriptsize
\begin{tabular}{llrrrcr}
\toprule
第一阶段规格 & 权重 & 模型 & $IMS$系数 & 标准误 & $p$值 & ESS \\
\midrule
核心共同原因 & 未截尾 & G1 & -0.7168 & 0.2746 & 0.0091 & 7,001.86 \\
核心共同原因 & 未截尾 & G2 & -0.2271 & 0.1479 & 0.1247 & 7,001.86 \\
核心共同原因 & P0.5--P99.5 & G1 & -0.6452 & 0.2700 & 0.0169 & 8,390.22 \\
核心共同原因 & P0.5--P99.5 & G2 & -0.2458 & 0.1150 & 0.0325 & 8,390.22 \\
核心共同原因 & P1--P99 & G1 & -0.5910 & 0.2633 & 0.0248 & 8,611.55 \\
核心共同原因 & P1--P99 & G2 & -0.2272 & 0.1153 & 0.0489 & 8,611.55 \\
\midrule
完整共同原因 & 未截尾 & G1 & -0.7051 & 0.2728 & 0.0098 & 6,856.92 \\
完整共同原因 & 未截尾 & G2 & -0.2162 & 0.1480 & 0.1441 & 6,856.92 \\
完整共同原因 & P0.5--P99.5 & G1 & -0.6528 & 0.2707 & 0.0159 & 8,345.77 \\
完整共同原因 & P0.5--P99.5 & G2 & -0.2490 & 0.1147 & 0.0300 & 8,345.77 \\
完整共同原因 & P1--P99 & G1 & -0.5897 & 0.2632 & 0.0250 & 8,596.82 \\
完整共同原因 & P1--P99 & G2 & -0.2278 & 0.1151 & 0.0477 & 8,596.82 \\
\midrule
完整共同原因加审计选择 & 未截尾 & G1 & -0.7584 & 0.3066 & 0.0134 & 7,197.35 \\
完整共同原因加审计选择 & 未截尾 & G2 & -0.2580 & 0.1270 & 0.0421 & 7,197.35 \\
完整共同原因加审计选择 & P0.5--P99.5 & G1 & -0.6150 & 0.2731 & 0.0243 & 8,078.75 \\
完整共同原因加审计选择 & P0.5--P99.5 & G2 & -0.2345 & 0.1195 & 0.0497 & 8,078.75 \\
完整共同原因加审计选择 & P1--P99 & G1 & -0.5485 & 0.2635 & 0.0373 & 8,437.53 \\
完整共同原因加审计选择 & P1--P99 & G2 & -0.2339 & 0.1183 & 0.0481 & 8,437.53 \\
\bottomrule
\end{tabular}
\begin{tablenotes}\footnotesize
\item 注：所有模型包含年度和行业固定效应，标准误按企业聚类。G1 只在结果方程中加入 $IMS$；G2 进一步加入全部滞后共同原因变量。三组第一阶段规格的条件密度和权重独立估计。
\end{tablenotes}
\end{threeparttable}
\end{table}

各规格中 $IMS$ 系数均为负，说明较高行业专长强度与较低漂绿之间的方向关系具有一致性。核心规格中，G2 的 $IMS$ 系数从未截尾权重下的 $-0.2271$ 变为 P0.5--P99.5 截尾下的 $-0.2458$，再变为 P1--P99 截尾下的 $-0.2272$。仅加入治理变量的完整共同原因规格下，G2 系数分别为 $-0.2162$、$-0.2490$ 和 $-0.2278$；进一步加入审计选择变量后，G2 系数分别为 $-0.2580$、$-0.2345$ 和 $-0.2339$。这些变化体现了治理变量、审计选择变量和极端权重共同改变了条件处理密度、目标伪总体、有效信息量及结果方程精度。

因此，不能将某一个显著性最强的结果视为唯一结论。更稳妥的表述是：在处理前 Greenwashing、企业财务、治理和审计选择变量所定义的可观测共同原因条件下，连续 $IMS$ 与漂绿之间保持负向关系；该关系在权重截尾和增强结果方程下总体保持，但统计精度取决于共同支持和极端权重处理。

\section{结论与局限}

本 GPS-IPW 设计将第一阶段协变量选择建立在企业选择高行业专长审计师的经济机制上。滞后漂绿、财务状况、企业生命周期、所有权和治理变量刻画企业的审计需求、风险和披露基础；滞后 Big4、审计费、审计任期和换所变量直接刻画企业--审计师匹配和行业专长选择；年份和行业固定效应控制共同制度和市场环境。该逻辑与 Bia and Mattei (2012)、Kluve et al. (2012) 和 Ammar and Boughrara (2023) 对处理前共同原因的处理一致。

当前逐剂量平衡结果显示，GPS 加权后 $|SMD|<0.1$ 的组合比例从约 74.7--75.0\% 提高至 92.5--93.3\%，平均绝对 SMD 从 0.0825 降至 0.0429，但最大绝对 SMD 仍约为 0.1666。因此，GPS 权重显著改善而非完全消除了可观测协变量不平衡。正式因果解释还需要对残余不平衡、共同支持、未观测混淆以及行业内溢出进行进一步敏感性分析。

\begin{thebibliography}{99}

\bibitem[Imbens(2000)]{imbens2000}
Imbens, G. W. 2000. The role of the propensity score in estimating dose-response functions. \textit{Biometrika} 87 (3): 706--710. \url{https://doi.org/10.1093/biomet/87.3.706}.

\bibitem[Hirano and Imbens(2004)]{hiranoimbens2004}
Hirano, K., and G. W. Imbens. 2004. The propensity score with continuous treatments. In A. Gelman and X.-L. Meng, eds., \textit{Applied Bayesian Modeling and Causal Inference from Incomplete-Data Perspectives}, 73--84. Wiley. DOI: \url{https://doi.org/10.1002/0470090456.ch7}.

\bibitem[Bia and Mattei(2012)]{biamattei2012}
Bia, M., and A. Mattei. 2012. Assessing the effect of the amount of financial aids to Piedmont firms using the generalized propensity score. \textit{Statistical Methods \& Applications} 21 (4): 485--516. \url{https://doi.org/10.1007/s10260-012-0193-4}.

\bibitem[Kluve et~al.(2012)]{kluve2012}
Kluve, J., H. Schneider, A. Uhlendorff, and Z. Zhao. 2012. Evaluating continuous training programmes by using the generalized propensity score. \textit{Journal of the Royal Statistical Society: Series A} 175 (2): 587--617. \url{https://doi.org/10.1111/j.1467-985X.2011.01000.x}.

\bibitem[Ammar and Boughrara(2023)]{ammarboughrara2023}
Ammar, N., and A. Boughrara. 2023. Noninterest income mix and bank performance: Generalized propensity score fractional dose--response function evidence. \textit{Global Finance Journal} 56: 100771. \url{https://doi.org/10.1016/j.gfj.2022.100771}.

\bibitem[Robins et~al.(2000)]{robinsetal2000}
Robins, J. M., M. A. Hern\'an, and B. Brumback. 2000. Marginal structural models and causal inference in epidemiology. \textit{Epidemiology} 11 (5): 550--560. \url{https://doi.org/10.1097/00001648-200009000-00011}.

\end{thebibliography}

\end{document}
